\documentclass[10pt]{article}
\pdfmapfile{=lm.map}
\pdfmapfile{+cm.map}
\pdfmapfile{+symbols.map}
\usepackage[T1]{fontenc}
\usepackage{lmodern,amsmath,amssymb,amsthm,booktabs,array,geometry,microtype,xcolor,hyperref}
\geometry{margin=1in}
\hypersetup{pdftitle={Counterexamples to Yust's Pitch-Space and Pitch-Class Folding Conjectures},pdfauthor={GPT-6 Astra},colorlinks=true,linkcolor=blue!45!black,urlcolor=blue!45!black,citecolor=blue!45!black}
\newtheorem{theorem}{Theorem}
\newtheorem{lemma}[theorem]{Lemma}
\newtheorem{definition}[theorem]{Definition}
\theoremstyle{remark}
\newtheorem{remark}[theorem]{Remark}
\newcommand{\Z}{\mathbb Z}
\newcommand{\N}{\mathbb N}
\newcommand{\ti}{\sim_{\mathrm{TI}}}
\newcommand{\multdiff}{\mathbin{\backslash_{\!m}}}
\newcommand{\card}[1]{\lvert #1\rvert}
\title{Counterexamples to Yust's Pitch-Space\\and Pitch-Class Folding Conjectures}
\author{GPT-6 Astra\\\small Project AI author label}
\date{14 October 2026}
\begin{document}
\maketitle
\begin{abstract}
We give explicit counterexamples to both arbitrary-cardinality folding
characterizations proposed by Yust in 2022. A seven-note integer pitch pair
shares two different six-note types and twenty of twenty-one interval
occurrences, but its first set has no inversion-symmetric five-note core and
admits no folding. A five-note pair in twelve-tone pitch-class space shares
two different four-note types and nine of ten interval-class occurrences,
but admits neither folding nor T-shift. Both concrete counterexamples have
kernel-checked Lean~4 certificates, including arbitrary source
transposition/inversion representatives and the integer-multiple form of
T-shift. We also prove an infinite pitch-space family in ordinary mathematics.
The formal certificates and their correspondence with the source statements
passed independent review and an isolated rebuild. Established novelty or
publication priority is not claimed.
\end{abstract}

\section{Statements and conventions}
Yust's Conjecture~1 characterizes folding of equal-cardinality pitch sets by
two shared codimension-one subset types and all but one pairwise interval
occurrence. Conjecture~2 proposes the analogous pitch-class characterization,
allowing either folding or T-shift \cite[p.~8]{Yust}. The preceding discussion
and Proposition~4 supply the relevant conventions: subset types are taken up
to transposition and inversion, interval content retains multiplicity, and an
$n$-note folding requires an inversion-symmetric $(n-2)$-note core
\cite[pp.~4--7]{Yust}.

We work in $G=\Z$ or $G=\Z/12\Z$ and write
\[
 T_s(x)=s+x,\qquad I_s(x)=s-x.
\]
For finite subsets $S,T\subset G$, write $S\ti T$ when $T=T_s(S)$ or
$T=I_s(S)$ for some $s$. In integer pitch space, let
\[
 \Delta(S)=\mathop{\biguplus}_{\substack{u,v\in S\\u<v}}\{v-u\}
\]
be the multiset of unordered pairwise intervals. In $\Z/12\Z$, use instead
the smaller modular distance in $\{1,\ldots,6\}$; its multiplicity vector is
denoted $\nu(S)=(\nu_1(S),\ldots,\nu_6(S))$. Thus the number of shared
interval occurrences is the cardinality of the multiset intersection, or,
in the pitch-class case, $\sum_{j=1}^6\min(\nu_j(S),\nu_j(T))$.

\begin{definition}[Operations used in the source]
Choose distinct $a,b\in S$ and put $C=S\setminus\{a,b\}$.
A folding with index $s$ requires $I_s(C)=C$ and replaces $S$ by
\[
 (S\setminus\{a\})\cup\{s-a\}.
\]
A T-shift requires $T_x(C)=C$ and, for some $k\in\Z$, replaces $S$ by
\[
 (S\setminus\{a\})\cup\{a+kx\}.
\]
The output must retain the original cardinality. When comparing types, source
and target representatives may be changed by transposition or inversion.
\end{definition}

Each operation is a \emph{single} folding or shift, not a sequence of such
operations. Each counterexample below has distinct notes, inequivalent whole
sets, two \emph{different} shared subset types, and exactly
$\binom n2-1$ shared interval occurrences. Accordingly, the failures do not
depend on repeated-note degeneracies, counting one subset type twice, or
interpreting ``share'' as ``at least'' rather than ``exactly.''

\section{A seven-note pitch counterexample}
\begin{theorem}\label{thm:pitch}
The integer pitch sets
\[
 A=\{0,1,2,3,6,7,9\},\qquad B=\{0,1,2,3,6,8,9\}
\]
share two different six-note types and exactly twenty pairwise interval
occurrences. They are not equivalent under transposition or inversion, and no
representative of $A$ admits a folding to any target. Thus the forward
implication in Yust's Conjecture~1 is false.
\end{theorem}
\begin{proof}
The first common six-set is
\[
 A\setminus\{7\}=B\setminus\{8\}=\{0,1,2,3,6,9\}.
\]
For the second, inversion $I_9$ gives
\[
 I_9(A\setminus\{2\})
 =I_9(\{0,1,3,6,7,9\})
 =\{0,2,3,6,8,9\}=B\setminus\{1\}.
\]
The two source six-sets are inequivalent: their common endpoints $0,9$
force any transposition between them to be $T_0$ and any inversion to be
$I_9$, and neither maps one to the other.

Their interval multiplicities are
\[
\begin{array}{c|rrrrrrrrr}
\text{length}&1&2&3&4&5&6&7&8&9\\\hline
\Delta(A)&4&3&3&2&2&3&2&1&1\\
\Delta(B)&4&3&3&1&2&3&2&2&1
\end{array}
\]
so, writing $\multdiff$ for multiset difference,
\[
 \Delta(A)\multdiff\Delta(B)=\{4\},\qquad
 \Delta(B)\multdiff\Delta(A)=\{8\}.
\]
The shared multiset has cardinality $20=\binom72-1$. Different interval
multisets also exclude $A\ti B$.

An odd finite reflection-invariant set contains the reflection center.
A symmetric five-note subset of $A$ would therefore require a center in $A$
and two distinct reflected pairs. The available pairs are
\[
\begin{array}{c|l}
\text{center}&\text{reflected pairs in }A\\\hline
0,9&\text{none}\\
1&\{0,2\}\\
2&\{1,3\}\\
3&\{0,6\}\\
6&\{3,9\}\\
7&\text{none}.
\end{array}
\]
No center supports two pairs, so no symmetric five-note core exists.
Consequently $A$ admits no seven-note folding. Transposition and inversion
preserve the existence and cardinality of symmetric subsets, so the same
obstruction holds for every representative of $A$.
\end{proof}

\section{An infinite pitch-space family}
\begin{theorem}\label{thm:family}
For every integer $t\ge6$, the sets
\[
\begin{split}
 A_t&=\{0,1,2,3,t,t+1,t+3\},\\
 B_t&=\{0,1,2,3,t,t+2,t+3\}
\end{split}
\]
are counterexamples of the same kind as in Theorem~\ref{thm:pitch}.
The exchanged interval occurrences are $t-2$ and $t+2$.
\end{theorem}
\begin{proof}
Their common six-set is
\[
 C_t=A_t\setminus\{t+1\}=B_t\setminus\{t+2\}
     =\{0,1,2,3,t,t+3\}.
\]
Also,
\[
 I_{t+3}(A_t\setminus\{2\})=B_t\setminus\{1\}.
\]
The source six-sets have endpoints $0,t+3$. Direct comparison excludes both
$T_0$ and $I_{t+3}$ as maps between them, so the shared types are different.

The distances from the exceptional point to $C_t$ are respectively
\[
\begin{split}
 \{|t+1-c|:c\in C_t\}&=\{t+1,t,t-1,t-2,1,2\},\\
 \{|t+2-c|:c\in C_t\}&=\{t+2,t+1,t,t-1,2,1\},
\end{split}
\]
where both sides are multisets. All distances within $C_t$ remain unchanged.
Hence the full multisets exchange precisely $t-2$ for $t+2$, have twenty
shared occurrences, and are different. Thus $A_t\not\ti B_t$.

Suppose a symmetric five-note subset exists. Its center is a member of
$A_t$. A center in the lower block $\{0,1,2,3\}$ supports at most one
reflected pair wholly in that block. A pair involving the upper block has
sum at least $t\ge6$, whereas twice the center is at most $6$. Equality
is possible only for $t=6$, center $3$, and pair $\{0,6\}$. That center
supports no reflected pair wholly in the lower block.

A center in the upper block $\{t,t+1,t+3\}$ has twice its value at least
$2t$. A pair involving the lower block has sum at most $t+6$. Equality is
possible only for $t=6$, center $6$, and pair $\{3,9\}$. That center has no
pair wholly in the upper block. The upper block itself contains no
three-term arithmetic progression. In all cases the center supports at
most one pair, contradicting the required five-note symmetry.
\end{proof}

\begin{remark}[Formalization boundary]
Theorem~\ref{thm:family} is an ordinary mathematical proof. The current Lean
certificate proves its concrete $t=6$ instance and the pitch-class example
below. It does not assert a formal theorem for the parameterized family.
Even-cardinality symmetric subsets are not excluded by the argument; the
obstruction is specifically to a five-note core.
\end{remark}

\section{A five-note pitch-class counterexample}
\begin{theorem}\label{thm:pc}
In $\Z/12\Z$, let
\[
 P=\{0,1,3,4,6\},\qquad Q=\{0,1,3,4,7\}.
\]
These sets share two different four-note types and exactly nine interval-class
occurrences. They are not equivalent under transposition or inversion and are related
by neither folding nor T-shift, even after arbitrary representative changes.
Thus the forward implication in Yust's Conjecture~2 is false.
\end{theorem}
\begin{proof}
The first common four-set is
\[
 P\setminus\{6\}=Q\setminus\{7\}=\{0,1,3,4\}.
\]
For the second,
\[
 I_7(P\setminus\{1\})=I_7(\{0,3,4,6\})
 =\{1,3,4,7\}=Q\setminus\{0\}.
\]
These source subset types differ because $\{0,1,3,4\}$ contains two
semitone pairs, while $\{0,3,4,6\}$ contains one.

Their interval-class vectors are
\[
 \nu(P)=(2,2,3,1,1,1),\qquad
 \nu(Q)=(2,1,3,2,1,1).
\]
They exchange one ic$2$ for one ic$4$, so their overlap is
$9=\binom52-1$ and $P\not\ti Q$.

The only inversion-symmetric three-note subset of $P$ is $\{0,3,6\}$,
with index $6$. To check this directly, the center of such an odd subset
must be a member $r\in P$ and its other two elements must sum to $2r$
modulo $12$. Among the possible centers $0,1,3,4,6$, only $r=3$
has such a pair, namely $\{0,6\}$. Folding either of the remaining notes
therefore gives
\[
 1\longmapsto5:\quad\{0,3,4,5,6\},\qquad
 4\longmapsto2:\quad\{0,1,2,3,6\}.
\]
Both outputs have vector $(3,2,2,1,1,1)$, unlike $Q$, so neither reaches
its type.

A nonidentity translation stabilizing a three-element subset has order
three: its equal-sized translation orbits partition that subset. Such a
subset is necessarily an augmented triad $\{r,r+4,r+8\}$. The set $P$
contains no augmented triad, so none of its three-note cores has a
nonidentity translational symmetry. The identity shift leaves $P$
unchanged, which is not $Q$'s type.

Finally, stabilizing translations form a subgroup. Every integer multiple
of a stabilizing translation still stabilizes the same core, so testing all
stabilizing translations includes the complete definition of T-shift.
Transposition/inversion changes of source representatives conjugate these
operations to operations of the same form and preserve core symmetries.
The obstruction consequently applies to every representative of $P$.
\end{proof}

\section{Formal certificate and independent verification}
The file \texttt{formal/Counterexamples.lean} models pitch sets by
\texttt{Finset} $\Z$ and pitch-class sets by \texttt{Finset (Fin 12)}.
It defines unordered pairwise interval multisets directly and retains their
multiplicities. The main exported theorems are
\begin{center}
\texttt{YustFolding.conjecture\_one\_false},\qquad
\texttt{YustFolding.conjecture\_two\_false}.
\end{center}
The formal conjecture statements include $n\ge3$, inequivalent whole sets,
two distinct shared deletion types, exact interval overlap, arbitrary source
representatives, and the printed integer-multiple T-shift predicate. These
are stricter content-side hypotheses than the printed conjectures explicitly
state; the counterexamples satisfy them all.

\paragraph{Unbounded pitch space.}
If a nonempty subset $C\subset[0,9]\cap\Z$ satisfies $I_s(C)=C$, then for
$x\in C$ we have $s-x\in C$, and hence $s=x+(s-x)\in[0,18]$. The Lean
proof establishes this bound before checking all possible integer indices.
It does not replace an unbounded pitch universe by a bounded one. Pullback
lemmas transport the obstruction to every transposition and inversion of
$A$: the pulled-back inversion indices are respectively $s-2u$ and
$2u-s$ for a source change by $T_u$ or $I_u$.

\paragraph{Pitch-class representatives and shifts.}
Every transposed and inverted source is checked; the operation predicates
already permit all target T/I representatives. The stabilizing translations
are explicitly packaged as an additive subgroup, and closure under integer
multiplication proves that the exact T-shift relation of Definition~4 is
included in the checked relation.

\paragraph{Kernel and audit.}
All finite certificates use ordinary kernel-checked \texttt{decide}. There
is no admitted proof, new axiom, \texttt{native\_decide}, or external solver
oracle. The exported theorem dependencies are only the standard foundations
\texttt{propext}, \texttt{Classical.choice}, and \texttt{Quot.sound}.
Python enumeration supported discovery and independent numerical checks but
is not imported as a proof dependency.

The frozen source passed an independent semantic review and a fresh isolated
Lean compilation, with exit status $0$ in $74.84$ seconds. No project object
files were reused. The audit checked the exact exported statements and the
source-to-formal bridges, including multiplicities, distinct subset types,
representatives, unbounded axes, and integer-multiple shifts. It found no
remaining semantic or proof gap. Receipts and exported statements are
included in the accompanying archive. The infinite family was independently
reviewed as an ordinary mathematical proof, with its formalization boundary
retained.

\section{Literature check and limitations}
The author-hosted revised paper and the author's publication list
\cite{Yust,YustList}, containing entries through 2026, were checked on
14 October 2026. Targeted exact-title searches combined with ``proof'' and
``counterexample,'' and searches for the two named conjectures, found no
later primary-source resolution. The Springer chapter landing page and the
two linked presentation files were not readable through the available
retrieval route. This is a bounded literature check, not proof of originality
or absence of prior work.

The results refute the arbitrary-cardinality characterizations. They do not
refute the proved pitch-space tetrachordal proposition, classify all
exceptional pairs, establish the smallest pitch-space counterexample, or
prove a repaired characterization. No publication-priority claim is made.

\appendix
\section{Reproducibility}
The toolchain is Lean~4.32.2, with mathlib pinned to
\begin{center}
\texttt{905b95818eb32af7874a58b427f50c1711a5e96c}.
\end{center}
The SHA256 of the audited one-module formal source is
\begin{center}\small\ttfamily
a884ebb6566c9920f3890c3b9b331a84c7e4a168404482847992a59b92217ace.
\end{center}
The archive includes \texttt{lean-toolchain}, \texttt{lakefile.toml}, the full
formal source, build logs, independent receipts, enumeration code and results,
this \LaTeX\ source, and the compiled paper. On a machine with Elan/Lake,
the standard build sequence is
\begin{verbatim}
lake update
lake exe cache get
lake build
\end{verbatim}
The supplied \texttt{formal/check.sh} additionally records the command used
with the existing research toolchain. The paper can be rebuilt with
\texttt{build-paper.sh}; its local writable TeX format avoids changing system
configuration.

\begin{thebibliography}{9}
\bibitem{Yust}
Jason Yust.
\newblock Tetrachordal Folding Operations.
\newblock In \emph{Mathematics and Computation in Music}, MCM 2022,
pp.~292--303. Springer, 2022.
\newblock DOI: \href{https://doi.org/10.1007/978-3-031-07015-0_24}{10.1007/978-3-031-07015-0\_24}.
\newblock \href{https://sites.bu.edu/jyust/files/2022/09/tetrachordalFoldingRev.pdf}{Author-hosted revised paper}, especially pp.~4--8.

\bibitem{YustList}
Jason Yust.
\newblock \href{https://sites.bu.edu/jyust/Publications/}{Publications}.
\newblock Boston University. Accessed 14 October 2026.
\end{thebibliography}
\end{document}
